%%%%%%%%%%%%%%%%%%%%%%%%%%%%%%%%%%%%%%%%%%%%%%%%%%%%%%%%%%%%%%
%%%% MA-0661 Algebra Abstracta II %%%%%
%%%%%%%%%%%%%%%%%%%%%%%%%%%%%%%%%%%%%%%%%%%%%%%%%%%%%%%%%%%%%%
\newpage
\subsection*{MA-0661 Algebra Abstracta II}
\addcontentsline{toc}{subsection}{MA-0661 Algebra Abstracta II}

\hypertarget{MA0661}{}
\begin{refsection}
\begin{table}[H]
\centering{
\begin{tabular}{|ll|}
\hline
%\multicolumn{2}{|c|}{\textbf{Nivel de virtualidad del curso:} Virtual}\\ \hline
\textbf{Año:} III  & \textbf{Requisitos:} MA-0561 Algebra Abstracta I \\ 
\textbf{Ciclo:} VI  & \textbf{Correquisitos:} Ninguno \\ 
\textbf{Tipo de curso:} Teórico & \textbf{Horas presenciales por semana:}   \\ 
\textbf{Créditos:} 4 & \textbf{Horas de trabajo independiente por semana:}  \\ \hline
\end{tabular}
}
\end{table}



\subsubsection*{Descripción del curso}

Este segundo curso de Álgebra Abstracta profundiza en el estudio de estructuras algebraicas fundamentales, con énfasis en la teoría de grupos, cuerpos y teorías asociadas como la de Galois. Se inicia con el análisis de la estructura interna de los grupos mediante acciones de grupo, teoremas de Sylow y aplicaciones, para luego introducir la teoría de cuerpos y extensiones, desarrollando conceptos clave como grado, elementos algebraicos y trascendentales.

Posteriormente, el curso aborda los automorfismos y el grupo de Galois, así como la teoría de extensiones normales y separables, culminando con el Teorema Fundamental de la Teoría de Galois y el estudio de cuerpos finitos. Se privilegia la comprensión estructural y la conexión entre objetos algebraicos, fomentando la capacidad de abstracción, formalización y razonamiento lógico-matemático.

\subsubsection*{Objetivos}

\textbf{General:} Analizar y aplicar los conceptos fundamentales de la teoría de grupos, cuerpos y extensiones de cuerpos, para comprender la estructura de las extensiones algebraicas y su relación con los grupos de automorfismos, mediante el uso de resultados clásicos de álgebra. \\

\textbf{Específicos:}

Al final del curso la persona estudiante debe ser capaz de:

\begin{enumerate}
\item Aplicar las definiciones fundamentales de grupos, cuerpos, extensiones, automorfismos y grupo de Galois.

\item Explicar la acción de un grupo sobre un conjunto, incluyendo órbitas, estabilizadores y sus propiedades asociadas.

\item Aplicar los teoremas de Sylow y los resultados básicos de la teoría de grupos para analizar la estructura de grupos finitos.

\item Analizar y aplicar conceptos de teoría de cuerpos y extensiones, incluyendo grado, elementos algebraicos y trascendentales, polinomio minimal y extensiones simples.

\item Analizar y relacionar las propiedades de normalidad y separabilidad de extensiones con la estructura del grupo de Galois y sus acciones.

\item Evaluar y construir estructuras algebraicas avanzadas (extensiones de Galois y cuerpos finitos) y resolver problemas que integren teoría de grupos y cuerpos mediante el uso del Teorema Fundamental de Galois.
\end{enumerate}

\subsubsection*{Contenidos}

\begin{enumerate}
    
\item \textbf{Teoremas de estructura de grupos (4 semanas):} 
\begin{enumerate}
\item Productos directos, grupos simétricos, alternantes y dihédricos.
    \item La acción de un grupo sobre un conjunto. Ejemplos.
    \item Órbitas, estabilizador, el teorema de la órbita y el estabilizador.
    \item Los teoremas de Sylow.
    \item Aplicaciones de los teoremas de Sylow.
\end{enumerate}

\item \textbf{Cuerpos (2 semanas):} 
\begin{enumerate}
    \item Cuerpos. Ejemplos básicos ($\Q, \R, \C, \mathbb{F}_p$). La característica de un cuerpo.
    \item Extensiones de cuerpos. El grado de una extensión. El Teorema de la Torre.
    \item Generadores de cuerpos y su caracterización. El compuesto.
    \item Elementos algebraicos y trascendentales en extensiones. El polinomio minimal.
    \item Extensiones simples y algebraicas.
\end{enumerate}

\item \textbf{Automorfismos (2 semanas):} 
\begin{enumerate}
    \item Automorfismos. $F$-homomorfismos.
    \item El grupo de Galois de una extensión. Su acción sobre las raíces de un polinomio sobre el cuerpo base.
  %  \item Cuerpos intermedios y cuerpos fijos de subgrupos del grupo de Galois y su correspondencia.
   % \item 
\end{enumerate}

\item \textbf{Extensiones normales (2 semanas)}:
\begin{enumerate}
    \item Definición, ejemplos. 
    \item Cuerpo de escisión. 
    \item Clausura algebraica. 
    \item El teorema de extensión de isomorfismos. 
    \item Condiciones equivalentes a la normalidad. 
\end{enumerate}

\item \textbf{Extensiones separables e inseparables (2 semanas):}
\begin{enumerate}
    \item Definición y ejemplos. 
    \item Propiedades básicas de la separabilidad. 
 \item Caracterización de extensiones de Galois en términos de separabilidad y normalidad. 
 \item Cuerpos perfectos. 
 \item Extensiones puramente inseparables (opcional). 
\end{enumerate}

\item \textbf{Teorema Fundamental de la Teoría de Galois (2 semanas):}

\begin{enumerate}
    \item Teorema fundamental de Galois.
    \item Ejemplos concretos no triviales. 
    \item Teorema del elemento primitivo. 
    \item La clausura normal de un extensión. 
\end{enumerate}


\item \textbf{Cuerpos finitos (1 semanas):}
\begin{enumerate}
       \item Descripción de los subgrupos finitos del grupo multiplicativo de un cuerpo como un grupo cíclico. 
    \item Teorema de estructura de los cuerpos finitos: Caracterización de extensiones finitas como cuerpos de escisión y estructura de su grupo de Galois, como grupo cíclico generado por el automorfismo de Frobenius. 
    \item Clasificación de extensiones de cuerpos finitos. 
    \item Construcción explícita de cuerpos finitos. 
\end{enumerate}

\end{enumerate}
\subsubsection*{Metodología}

\noindent \textbf{Lineamientos metodológicos:} La persona docente a cargo de este curso debe respetar las siguientes pautas:
\begin{enumerate}

    \item Fomentar una visión integral de las matemáticas y de cómo se enlazan las diferentes áreas a través del álgebra. Por ejemplo sus relaciones con geometría algebraica, topología algebraica, teoría de números, criptografía, lógica, etc.
\end{enumerate}

\textbf{Sugerencias metodológicas:} Adicionalmente, se tienen las siguientes recomendaciones para la persona docente a cargo de este curso:
\begin{enumerate}

    \item Utilizar el recurso tecnológico para reforzar distintos conceptos estudiados en el curso por medio de la visualización, cálculo y experimentación.
    \item Aprovechar momentos adecuados del curso para motivar el estudio por medio de reseñas acerca del desarrollo histórico del álgebra y su importancia en aplicaciones dentro y fuera de la matemática
    \item En la evaluación, se sugiere utilizar estrategias que promuevan el trabajo constante de parte de las personas estudiantes a lo largo del ciclo lectivo. Por ejemplo, esto se puede hacer por medio de tareas o quizzes.
    \item Para fomentar el desarrollo de las habilidades investigativas en el estudiantado, se sugiere la implementación de proyectos de investigación y participación en coloquios o seminarios de investigación.
    \item Promover la comunicación clara y rigurosa de argumentos e ideas matemáticas de manera oral y escrita.
\end{enumerate}




\nocite{Alu09}
\nocite{Alu21}
\nocite{CE20}
\nocite{Mor96}
\nocite{New12}
\nocite{Edw84}
\nocite{Sti10}
\nocite{Kle07}
\nocite{Kna16a}
\nocite{Kna16b}
\nocite{Lan02}
\nocite{Lan05}

\printbibliography[heading=subbibliography]
\end{refsection}
